\section{Discussion}
\label{sec:discussion}

\textbf{What the debate structure contributes.}
The strongest evidence for \abbrev{} is not only that its final scores are higher, but that the improvement pattern matches the intended mechanism. A compute-matched zero-shot baseline evaluates the same 80 candidates without iterative feedback and remains 18--35\% below \abbrev{}, while single-agent and thesis-only ablations plateau earlier. The cross-over study in Table~\ref{tab:design_reward_ablation} shows why this matters for co-design: in Hopper, the \abbrev{} morphology alone underperforms the default body, yet the co-designed reward recovers a 50\% gain. This is the failure mode that morphology-only or reward-only search cannot address cleanly. The role-separated debate gives the design agent a structured way to propose body changes, gives the control agent an explicit channel to flag controllability issues before training, and lets the archive retain empirical winners when individual proposals fail.

\textbf{What the judges do and do not do.}
The judge panel should not be interpreted as an LLM reward model. Judges do not select winners, assign scalar fitness, or replace simulation. They only translate measured behavior into textual critiques for the next round, while the archive is ranked by the external score $S$. This separation is important because it prevents persuasive but ungrounded LLM explanations from determining the final output. The evidence is consistent with judges acting as a directional search aid: speed and stability critiques predict leg and torso edits, and judge-conditioned runs make smaller, more targeted parameter changes than no-judge runs (Section~\ref{sec:ablation_studies}; Appendix~\ref{app:critique-analysis}). Thus the judges improve the proposal distribution rather than changing the objective.

\textbf{When the framework transfers.}
Applying \abbrev{} to a new domain requires more than an LLM and a simulator: the domain must provide a constrained design parameterization, an executable reward template, and an external task score $S$ for archive selection. The third requirement is the most important. LLMs can propose reward code, but they should not define the metric that declares success. Under this contract, debate acts as a proposal mechanism for coupled design--reward artifacts, while simulator-measured performance remains the arbiter. Domains without repeatable simulation or task-level success metrics would require a different evaluation protocol and should be treated as new experiments rather than direct transfers.
